


% \comment[Oreofe]{We might need more than just saying \textit{...a size comparable to prior package-ecosystem measurement studies}; we can add the C.I. and MoE}
% \comment[Oreofe]{0) Dataset (I'm thinking we should include how we sourced our dataset and mined it. We did some work, especially for PyPI and npm, since there wasn't an attestation dataset. And we can state our adoption rates that we found. Are the mined datasets also  a contribution?}

\section{Artifact Verifiability Methodology}
\label{sec:methodology}

This section describes our methodology for answering the research questions in \cref{sec:knowledge-gaps}.

\subsection{Experimental Setup}
\label{subsec:experimental-setup}

\subsubsection{Package Ecosystem Selection}
We study four source-oriented ecosystems: Crates.io, npm, PyPI, and RubyGems. These are among the largest programming-language ecosystems and frequent targets
of supply-chain attacks~\cite{ohm_backstabbers_2020}; they also match the source-oriented ecosystems studied by recent reproducible-build work~\cite{benedetti_empirical_2025} and supported by OSS Rebuild~\cite{google_oss_rebuild}. 
% We study four source-oriented package ecosystems: Crates.io, npm, PyPI, and RubyGems.
% These ecosystems are among the largest programming-language package ecosystems and are frequent targets of software supply chain attacks~\cite{ohm_backstabbers_2020}.
% They also match the source-oriented ecosystems studied by recent reproducible-build work~\cite{benedetti_empirical_2025} and supported by Google's OSS Rebuild~\cite{google_oss_rebuild}.
% \Cref{xx} discusses the generalizability of our findings to ecosystems outside this scope.

% \begin{table}
% \centering
% \caption{Dataset summary.}
% \label{tab:dataset-details}
% \begin{tabular}{lrrrr}
% \toprule
% Ecosystem & Packages & Valid repos & Attested & Sampled artifacts \\
% \midrule
% npm       & \mock{XX} & \mock{XX} & \mock{XX} & \mock{XX} \\
% PyPI      & \mock{XX} & \mock{XX} & \mock{XX} & \mock{XX} \\
% crates.io & \mock{XX} & \mock{XX} & \mock{XX} & \mock{XX} \\
% RubyGems  & \mock{XX} & \mock{XX} & \mock{XX} & \mock{XX} \\
% \bottomrule
% \end{tabular}
% \end{table}

% \begin{table}[t]
% \centering
% \caption{Dataset summary. }
% \label{tab:dataset-details}
% \begin{tabular}{lrrrr}
% \toprule
% Ecosystem & Packages & Valid repos & Attested & Sampled artifacts \\
% \midrule
% npm       & 9{,}000 & 8{,}964 & 4{,}506 & 7{,}607 \\
% PyPI      & 9{,}000 & 8{,}943 & 4{,}502 & 7{,}793 \\
% crates.io & 9{,}706\makebox[0pt][l]{\textsuperscript{*}} & 9{,}676 & 5{,}206 & 9{,}706 \\
% RubyGems  & 6{,}299\makebox[0pt][l]{\textsuperscript{*}} & 6{,}299 & 1{,}799 & 6{,}240 \\
% \bottomrule
% \end{tabular}
%   \begin{tablenotes}[para]
%     \small \item[*] Number of packages differ because  \mock{...}
%   \end{tablenotes}
% \end{table}

\subsubsection{Dataset Construction and Sampling}
\label{subsubsec:setup-artifact-selection}
For each ecosystem, we enumerate packages from the registry and build a package-level dataset of metadata, source-repository information, and provenance-attestation status. We recover repository URLs from registry fields such as project URL, homepage, and description, following prior
work~\cite{gao_pyradar_2024}, and query each ecosystem's provenance API for whether the latest artifact is attested.

Two properties of the dataset drive our sampling. Attestation adoption is low (2.0-6.6\%), so a random ecosystem-level sample would contain too few attested artifacts to evaluate attestations. Many packages lack valid repository URLs (14.7-42.4\%), which we report as an ecosystem-level source-recovery limitation. We exclude packages with no registry-linked public repository candidate from the rebuild experiment; however, \toolname may still fail source recovery for sampled packages when the linked repository is unavailable, invalid, lacks a matching package subpath, or no source commit can be recovered. We therefore sample within attested and non-attested strata separately, drawing up to 4,500 packages per stratum, above the sample needed for a 95\% confidence level and 1.5\% margin of error under the conservative infinite, population assumption of prior work~\cite{benedetti_empirical_2025}. For crates.io and RubyGems we take the full attested set---5{,}206 and 1{,}799 packages, respectively---since RubyGems has fewer than 4{,}500 attested packages in total, and for crates.io, we use all attested crates rather than subsampling. For each sampled package we select its most recent artifact as of May 2026, yielding 34{,}005 artifacts (\cref{tab:dataset-and-completion}); \Cref{fig:dataset-characteristics} shows the
download-rate distribution.

\begin{figure}[t]
\centering
\includegraphics[width=0.8\linewidth]{figures/downloads_distribution.pdf}
\caption{Download distribution of the sampled packages per ecosystem.}
\label{fig:dataset-characteristics}
\end{figure}

% For each ecosystem, we construct a package-level dataset containing package metadata, source repository information, and provenance-attestation status.
% We first enumerate packages from the corresponding registry.
% We then recover source repository URLs from registry metadata fields such as project URL, homepage, and description, following prior work~\cite{PyRadar}.
% For each package, we query the ecosystem-provided provenance-attestation APIs to identify whether the latest published artifact has an associated provenance attestation.


% Two properties of the resulting dataset shape our sampling strategy.
% First, provenance-attestation adoption is low across ecosystems, ranging from \mock{2.0\%} to \mock{6.6\%}.
% A purely random ecosystem-level sample would therefore contain too few attested artifacts to evaluate the effect of attestations.
% Second, many packages lack valid repository URLs, ranging from \mock{14.7\%} to \mock{42.4\%} across ecosystems.
% We report these packages as an ecosystem-level source-recovery limitation, but exclude non-attested packages without valid repository URLs from the rebuild experiment because they provide no public source candidate for an independent verifier.
% These exclusion rates are consistent with prior measurements of valid repository metadata in package ecosystems~\cite{}.

% We therefore use stratified sampling.
% For each ecosystem, we divide packages into attested and non-attested strata.
% \PA{For Crates.io, we use the full set because it was only 5209.}
% Within each stratum, we randomly sample up to 4,500 packages, which exceeds the 4,269 samples required to estimate proportions with a 95\% confidence level and a 1.5\% margin of error under the conservative infinite-population assumption used by prior work~\cite{benedetti_empirical_2025}.
% For each sampled package, we select the most recently published artifact as of May 2026, the time of our study.
% This yields \mock{XX} artifacts in total.
% \Cref{tab:dataset-details} summarizes the resulting dataset, and \cref{fig:dataset-characteristics} shows the distribution of package download rates.


\subsubsection{Experiment environment}
\label{subsubsec:rebuild-env}

We run all rebuilds and measurements on an IBM Cloud bare-metal server \iffalse (\textit{bx2-metal-96x384}\fi profile: 48 physical cores, 384\, GB RAM, and a 960\, GB local NVMe disk\iffalse provisioned in IBM Cloud's \texttt{eu-de} region\fi. The server runs Ubuntu Linux 24.04 LTS (x86-64)\iffalse (Noble Numbat) on the x86-64 architecture\fi. 


\subsection{RQ1: Proportion of Verified Software Packages}

RQ1 measures the proportion of packages in each ecosystem that can be independently verified under the models in \cref{sec:verifiability-model}. For each selected artifact, we run \toolname to recover rebuild metadata, reconstruct the build environment, rebuild the package twice, and compare the rebuilt artifacts against the published one. We run three experiments: (i) the non-attested ecosystem sample with heuristic metadata recovery, (ii) the attested
sample with heuristic source recovery, and (iii) the attested sample with attestation-based recovery. The latter two isolate the contribution of provenance attestations to verifiability.

We report two metrics per experiment: the \textit{completion rate} (the fraction of artifacts for which \toolname produces a comparison result, together with the reasons for pre-comparison failures), and the \textit{equivalence rate} (the fraction reaching each level). We consider an artifact \textit{verified} at L1, L2, or L3 (\cref{subsec:artifact-comparison-model}).

% RQ1 measures the proportion of packages in each ecosystem that can be independently verified under the verifiability models in \cref{sec:verifiability-model}.
% For each selected artifact, we run \toolname to recover rebuild metadata, reconstruct the build environment, rebuild the package twice, and compare the rebuilt artifacts against the published artifact.
% We run three experiments:
% (i) the non-attested ecosystem sample with heuristic-based metadata recovery,
% (ii) the attested sample with heuristic-based source recovery, and
% (iii) the attested sample with attestation-based metadata recovery.
% The second and third experiments allow us to isolate the contribution of provenance attestations to artifact verifiability.

% \begin{table}
%     \caption{Pipeline completion rates and reasons for pipeline failures.}
%     \centering
%     \begin{tabular}{lrrrr}
%         \toprule
%         Metric & Crates.io & NPM & PyPI & RubyGems \\
%         \midrule
%         Artifacts (dataset)         & 9{,}706 & 9{,}000 & 9{,}000 & 6{,}299 \\
%         \addlinespace
%         Completion rate             & 83.4\%  & 61.6\%  & 74.4\%  & 74.9\%  \\
%         \addlinespace
%         \multicolumn{5}{l}{\textit{Failure reasons (\% of dataset):}} \\
%         \;\;No source commit/repo   & 4.1\%   & 15.5\%  & 13.4\%  & 15.8\%  \\
%         \;\;Clone/download failure  & 9.1\%   & 2.5\%   & 0.4\%   & 2.1\%   \\
%         \;\;Build/install failure   & 3.4\%   & 20.4\%  & 11.2\%  & 7.3\%   \\
%         \;\;Timeout/other error     & 0.0\%   & 0.0\%   & 0.6\%   & 0.0\%   \\
%         \bottomrule
%     \end{tabular}
%     \label{tab:pipeline-completion}
% \end{table}


% \begin{table}
%     \centering
%     \begin{tabular}{ccccc}
%         Metric & Crates.io & NPM & PyPI & RubyGems \\
%          &  &  &  & \\
%          &  &  &  & \\
%          &  &  &  & \\
%          &  &  &  & \\
%          &  &  &  & \\
%     \end{tabular}
%     \caption{Pipeline completion rates and reasons for pipeline failures.}
%     \label{tab:placeholder}
% \end{table}

% For each experiment, we report two metrics.
% First, we measure the \textit{pipeline completion rate}: the proportion of artifacts for which \toolname completes and produces a comparison result.
% We also record the reasons for failures before comparison.
% Second, we measure the \textit{equivalence rate}: the proportion of artifacts assigned to each equivalence level.
% We consider an artifact \textit{verified} if \toolname classifies it as L1-, L2-, or L3-equivalent under the levels defined in \cref{subsec:artifact-comparison-model}.

\subsection{RQ2: Root Causes of Non-L1 Verifiability Failures}


For artifacts that did not reach L1, we qualitatively analyzed the observed differences and their causes. We restricted this to artifacts with provenance metadata, whose attestations or trusted-publishing records identify the source commit and CI/CD workflow, enabling inspection and partial reconstruction of the original build.

In the first stage, two authors independently open-coded 48 randomly sampled non-L1 artifacts: for each, they examined the \texttt{diffoscope}~\cite{diffoscope} report to identify observable differences, then inspected the source commit and GitHub
Actions workflow to infer the underlying source-state, build-process, or environment cause (30-60 minutes per artifact). They reconciled their findings into a shared taxonomy and codebook defining each symptom, its possible causes, and ecosystem-specific manifestations.

In the second stage, we used a frontier coding agent to scale the analysis, requiring experimental validation for each proposed cause: the agent corrected the hypothesized cause in a Dockerfile matching our rebuild environment and
checked whether the artifact's equivalence level improved, with a second agent reviewing each classification. Applied over three phases to saturation, this covered 134 further artifacts (35 from Crates/npm/PyPI and 29 from RubyGems). We report the resulting causes
and symptoms, grouped by the rebuild metadata they affect, with frequencies across ecosystems.

\begin{table}[t]
\caption{Dataset and pipeline completion per ecosystem. Completion rate and failure reasons are \% of sampled artifacts.}
\label{tab:dataset-and-completion}
\centering
\begin{tabular}{lrrrr}
\toprule
 & Crates.io & npm & PyPI & RubyGems \\
\midrule
Sampled: attested       & 5{,}206 & 4{,}500 & 4{,}500 & 1{,}799 \\
Sampled: non-attested   & 4{,}500 & 4{,}500 & 4{,}500 & 4{,}500 \\
Total sampled           & 9{,}706 & 9{,}000 & 9{,}000 & 6{,}299 \\
\midrule
% \vspace{1pt}
\;Completion rate         & 83.4\%  & 61.6\%  & 74.4\%  & 74.9\%  \\
\;\;\;\;No source/repo        & 4.1\%   & 15.5\%  & 13.4\%  & 15.8\%  \\
\;\;\;\;Clone/download        & 9.1\%   & 2.5\%   & 0.4\%   & 2.1\%   \\
\;\;\;\;Build/install         & 3.4\%   & 20.4\%  & 11.2\%  & 7.3\%   \\
\;\;\;\;Timeout/other         & 0.0\%   & 0.0\%   & 0.6\%   & 0.0\%   \\
\bottomrule
\end{tabular}
\end{table}

% For artifacts that did not achieve L1 equivalence, we qualitatively analyzed the observed differences and their root causes. We restricted this analysis to artifacts with provenance metadata because their attestations or trusted publishing metadata identify the source commit and CI/CD workflow used to produce the original artifact, enabling inspection and partial reconstruction of the original build process.

% Our analysis had two stages. First, two authors independently analyzed 48 randomly sampled non-L1 artifacts across the studied ecosystems. Using open coding, they examined each artifact’s \texttt{diffoscope} report to identify observable differences, then inspected the corresponding source commit and GitHub Actions workflow to infer the underlying source-state, build-process, or environment cause. 
% Each artifact took 30--60 minutes to analyze. 
% The authors periodically reconciled their findings into a shared taxonomy and codebook that defines each symptom, its possible root causes, and ecosystem-specific manifestations.

% Second, we used a frontier coding agent to scale the analysis. 
% This automation was necessary because manual root cause analysis was time-intensive and often required understanding package-specific build processes. 
% To mitigate agent non-determinism and errors, we required experimental validation for each proposed root cause: the agent created a Dockerfile matching our rebuild environment, modified it to correct the hypothesized cause, and checked whether the correction improved the artifact’s equivalence level.
% A second agent then reviewed the proposed classification and validation evidence. 

% We applied this process in three phases, refining the taxonomy and codebook after each phase until reaching saturation. 
% In total, this step covered 140 additional artifacts, sampled as 35 artifacts per ecosystem.
% We report the resulting root causes and failure symptoms, group them by the rebuild metadata they affect, and measure their frequencies across ecosystems.


% \subsection{RQ3: Heuristic vs Attestation-Based Source Recovery}

% RQ3 evaluates how source recovery strategies (\cref{subsubsec:stage-1-source-recovery}) affect artifact verifiability.
% This question is important because provenance attestations remain rare, making heuristic recovery necessary for most artifacts.
% We first measure heuristic accuracy for NPM and PyPI by comparing the repository URL, source commit, and package subpath recovered by heuristics against source metadata from attestations.
% We then compare the verifiability outcomes produced by heuristic-based and attestation-based recovery on the same provenance-attested packages.
% For artifacts verified at L1 or L2 by one method but not the other, we inspect the recovered metadata and build outputs to identify why one strategy did not yield an accurate source location.


\subsection{RQ3: Impact of Artifact Characteristics and Development Choices}

RQ3 tests whether six factors are associated with verifiability: artifact size, download rate, and age (artifact characteristics), and build tool, native code,
and monorepo layout (development choices). The last three capture development choices that may complicate rebuilding: monorepos can obscure package subpaths, native extensions can depend on platform-specific environments, and build tools may introduce different sources of nondeterminism.
We stratify continuous factors into quartiles and group categorical factors, comparing verifiability across bins.


% RQ3 examines whether artifact characteristics and development choices are associated with artifact verifiability.
% We consider six factors: artifact size, download rate, artifact age, monorepo use, native or platform-specific code, and build or packaging tool.
% The first three capture artifact characteristics that may influence reproducibility: larger artifacts provide more opportunities for divergence, while popular or recent artifacts may reflect current development practices.
% The last three capture development choices that may complicate rebuilding: monorepos can obscure package subpaths, native extensions can depend on platform-specific environments, and build tools may introduce different sources of nondeterminism.

% For continuous factors, we stratify artifacts into quartiles and compare verifiability rates and equivalence levels across quartiles. For categorical factors, we group artifacts by category and compare verifiability rates and equivalence levels across groups.


\subsection{RQ4: \toolname vs State-of-Practice Rebuild Platforms}
\label{subsec:method-rq5}

RQ4 compares \toolname with OSS-Rebuild, Google's open-source platform that infers build definitions, rebuilds artifacts, and produces rebuild attestations for popular packages. We select it as open source, targeting our ecosystems, and operated at scale. We compare on Crates.io, npm, and PyPI, excluding RubyGems, which OSS-Rebuild does not support.

We draw 600 artifacts: 200 per ecosystem, split into two strata of 100 by OSS-Rebuild coverage. The supported stratum is sampled from packages OSS-Rebuild already tracks; the unsupported stratum from each ecosystem's most popular
packages by download count, excluding any already supported. Both strata are thus widely-used software differing only in coverage, which separates how the platforms compare where OSS-Rebuild is tuned from how they compare on the popular
long tail it does not support. This frame is independent of the attestation-stratified dataset used in our other RQs.

We rebuild each artifact with both platforms from the same published version: OSS-Rebuild via its tool, which infers a build definition, builds hermetically, and reports whether the result matches; \toolname via our pipeline
(\cref{fig:method_overview}) on the same artifact and recovered source commit. Because OSS-Rebuild reports a binary match while \toolname reports graded levels (\cref{subsec:artifact-comparison-model}), we count an \toolname artifact as
verified at L1--L3, normalizing the same non-semantic differences OSS-Rebuild normalizes when deciding a match. We report, per ecosystem and stratum, each platform's verification rate and their agreement (verified by both, \toolname only, OSS-Rebuild only, or neither); for artifacts verified by exactly one, we analyze the design choices: source recovery, build-environment inference,
native-code handling, that explain the divergence.


% RQ4 compares \toolname with OSS-Rebuild, Google's open-source platform that infers build definitions, rebuilds artifacts, and produces rebuild attestations for popular packages. We select OSS-Rebuild because it is open source, targets our studied ecosystems, and is operated at scale by a major technology company. We compare on Crates.io, npm, and PyPI, excluding RubyGems, which OSS-Rebuild does not support.

% \paragraph{Sample.} We draw 600 artifacts: 200 per ecosystem split into two strata of 100 by OSS-Rebuild coverage. The supported stratum is sampled from packages OSS-Rebuild already tracks; the unsupported stratum from each ecosystem's most popular packages by download count, excluding any already supported. Both strata are thus widely-used software differing only in OSS-Rebuild coverage, separating how the platforms compare on packages OSS-Rebuild is tuned for from how they compare on the popular long tail it does not support. This frame is independent of the attestation-stratified dataset used in our other RQs.

% \paragraph{Procedure.} We rebuild each artifact with both platforms from the same published version. For OSS-Rebuild's tool, which infers a build definition, builds in a hermetic environment, and reports whether the rebuilt artifact matches. For \toolname we run our pipeline (??) on the same artifact and recovered source commit.

% \paragraph{Comparison.} OSS-Rebuild reports a binary outcome(match or no-match) , whereas \toolname produces graded equivalence levels (See \Cref{subsec:artifact-comparison-model}). For a fair comparison, we treat an \toolname artifact as verified at L1–L3, which normalizes the same classes of non-semantic difference OSS-Rebuild normalizes when deciding a match. We report, per ecosystem and stratum: (i) each platform's verification rate, and (ii) agreement—artifacts verified by both, \toolname only, OSS-Rebuild only, or neither. For artifacts verified by exactly one platform, we analyze the design choices (source recovery, build-environment inference, native-code handling) that explain the divergence.
